\documentclass[11pt]{article}

% ============================================================
% Packages
% ============================================================
\usepackage[utf8]{inputenc}
\usepackage[margin=1in]{geometry}
\usepackage{amsmath,amssymb,amsthm}
\usepackage{graphicx}
\usepackage{booktabs}
\usepackage{xcolor}
\usepackage{hyperref}
\usepackage{tikz}
\usepackage{enumitem}
\usepackage{pgfplots}

\pgfplotsset{compat=1.18}

% ============================================================
% Document Formatting
% ============================================================
\setlength{\parskip}{0.6em}
\setlength{\parindent}{0pt}

% ============================================================
% List Formatting
% ============================================================
\setlist[enumerate]{
    leftmargin=2em,
    itemsep=0.25em
}

% ============================================================
% Solution Environment
% ============================================================
\newenvironment{solution}
    {\par\noindent\textbf{Solution.}\par}
    {\par}

% ============================================================
% Link Styling
% ============================================================
\hypersetup{
    colorlinks=true,
    linkcolor=blue,
    urlcolor=cyan
}

% ============================================================
% Metadata
% ============================================================
\title{Ordinary Differential Equations}
\author{Chris Ritter}
\date{October 15, 2026}

% ============================================================
% Document
% ============================================================


\begin{document}

\maketitle
\tableofcontents
\newpage

% ============================================================
\section{Nonhomogeneous Linear Equations}
% ============================================================

A second-order nonhomogeneous linear differential equation has the form
\[
L[y]:=y''+p(t)y'+q(t)y=g(t).
\]
The associated homogeneous equation is \(L[y]=0\). If \(y_h\) is its
general solution and \(y_p\) is any one solution of \(L[y_p]=g\), then
\[
\boxed{y=y_h+y_p}
\]
is the general solution of the nonhomogeneous equation. Indeed,
linearity gives \(L[y_h+y_p]=L[y_h]+L[y_p]=g\); conversely, the
difference of any two solutions of \(L[y]=g\) solves \(L[y]=0\).

For constant coefficients, the homogeneous part comes from the
characteristic equation. The \textbf{method of undetermined
coefficients} then chooses a trial form for \(y_p\) whose derivatives
have the same types of terms as \(g\). It applies when \(g\) is a
finite sum of polynomials, exponentials, sines, cosines, and their
products. Once the trial is chosen, substitution determines its
unknown coefficients.

\subsection{Choosing a Trial Particular Solution}

For a polynomial \(P_n(t)\) of degree \(n\), use a general polynomial
\(Q_n(t)\) of the same degree. Typical trial forms are
\[
\begin{array}{c|c}
\text{Forcing term} & \text{Initial trial for }y_p\\ \hline
P_n(t) & Q_n(t)\\
e^{\alpha t}P_n(t) & e^{\alpha t}Q_n(t)\\
P_n(t)\cos\beta t\text{ or }P_n(t)\sin\beta t
  & Q_n(t)\cos\beta t+R_n(t)\sin\beta t\\
e^{\alpha t}P_n(t)\cos\beta t\text{ or }e^{\alpha t}P_n(t)\sin\beta t
  & e^{\alpha t}\bigl(Q_n(t)\cos\beta t+R_n(t)\sin\beta t\bigr)
\end{array}
\]
Include both sine and cosine terms because differentiation mixes them.
Likewise, include every polynomial power through degree \(n\).

\textbf{Resonance.} If any term in the trial solves the homogeneous
equation, the trial cannot produce the forcing term. Multiply the
entire trial by \(t^s\), where \(s\) is the multiplicity of the
corresponding characteristic root \(\alpha+i\beta\). Then substitute
and solve for the coefficients. For real coefficients, complex roots
come in conjugate pairs and share the same multiplicity.

\subsection{Example: Resonance with a Trigonometric Forcing}

Solve
\[
y''+4y=3\cos 2t.
\]

\begin{solution}
The characteristic equation \(r^2+4=0\) has roots \(\pm2i\), so
\[
y_h=C_1\cos2t+C_2\sin2t.
\]
The usual trial \(A\cos2t+B\sin2t\) is already homogeneous and
therefore gives zero when substituted into \(y''+4y\). The roots have
multiplicity one, so use
\[
y_p=t(A\sin2t+B\cos2t).
\]
Direct differentiation shows
\[
(D^2+4)(t\sin2t)=4\cos2t,
\qquad
(D^2+4)(t\cos2t)=-4\sin2t.
\]
Thus \(4A=3\) and \(B=0\). The general solution is
\[
\boxed{y=C_1\cos2t+C_2\sin2t+\frac34t\sin2t.}
\]
\end{solution}

\subsection{Example: A Repeated Root and Polynomial Forcing}

Solve
\[
y''-4y'+4y=t e^{2t}.
\]

\begin{solution}
The characteristic polynomial is
\(r^2-4r+4=(r-2)^2\). Hence
\[
y_h=(C_1+C_2t)e^{2t}.
\]
The forcing is an exponential times a degree-one polynomial, and
\(r=2\) is a double root. Multiply the ordinary trial by \(t^2\):
\[
y_p=t^2e^{2t}(At+B)=e^{2t}(At^3+Bt^2).
\]
For any smooth \(v\),
\[
(D-2)^2\bigl(e^{2t}v(t)\bigr)=e^{2t}v''(t).
\]
With \(v=At^3+Bt^2\), the equation becomes
\[
e^{2t}(6At+2B)=te^{2t}.
\]
Therefore \(A=1/6\) and \(B=0\), giving
\[
\boxed{y=e^{2t}\left(C_1+C_2t+\frac{t^3}{6}\right).}
\]
\end{solution}

\subsection{Superposition for Several Forcing Terms}

If \(g=g_1+\cdots+g_m\) and \(L[y_{p,j}]=g_j\), then
\[
L[y_{p,1}+\cdots+y_{p,m}]=g_1+\cdots+g_m.
\]
Find one particular solution for each forcing term and add them.
This works because \(L\) is linear. After adding \(y_h\), combine or
absorb any duplicate homogeneous terms into its constants.

\subsection{Example: Three Forcing Terms}

Find a particular solution of the equation shown in the class notes:
\[
y''-3y'-4y=3e^t+2\sin t-8e^t\cos2t.
\]

\begin{solution}
Write \(L=D^2-3D-4\). Split the right-hand side into three terms.

For \(3e^t\), try \(y_{p,1}=Ae^t\). Since
\(L[e^t]=(1-3-4)e^t=-6e^t\), we obtain
\[
y_{p,1}=-\frac12e^t.
\]

For \(2\sin t\), try
\(y_{p,2}=A\sin t+B\cos t\). Comparing coefficients gives
\[
-5A+3B=2,\qquad -3A-5B=0,
\]
so \(A=-5/17\) and \(B=3/17\). Thus
\[
y_{p,2}=-\frac5{17}\sin t+\frac3{17}\cos t.
\]

For \(-8e^t\cos2t\), try
\(y_{p,3}=e^t(A\cos2t+B\sin2t)\). Substitution gives
\[
-10A-2B=-8,\qquad 2A-10B=0,
\]
so \(A=10/13\) and \(B=2/13\). Therefore
\[
\boxed{y_p=-\frac12e^t-\frac5{17}\sin t+\frac3{17}\cos t
+\frac{10}{13}e^t\cos2t+\frac2{13}e^t\sin2t.}
\]
The homogeneous characteristic equation
\((r-4)(r+1)=0\) gives
\(y_h=C_1e^{4t}+C_2e^{-t}\). Hence the full general solution is
\[
\boxed{y=C_1e^{4t}+C_2e^{-t}+y_p.}
\]
\end{solution}

\subsection{When to Use a Different Method}

Undetermined coefficients relies on the forcing term staying in a
finite-dimensional family under differentiation. A forcing term such
as \(\tan t\), \(\csc t\), or \(\ln t\) generally requires another
method. \textbf{Variation of parameters} handles such terms when a
fundamental pair of homogeneous solutions is known.

\end{document}
